% greedy_baseline.tex
% -----------------------------------------------------------------------------
% Material for the B2 revision item ("why an evolutionary algorithm rather than
% something simpler?"). Included from main.tex at the end of Section VII via
%   % greedy_baseline.tex
% -----------------------------------------------------------------------------
% Material for the B2 revision item ("why an evolutionary algorithm rather than
% something simpler?"). Included from main.tex at the end of Section VII via
%   % greedy_baseline.tex
% -----------------------------------------------------------------------------
% Material for the B2 revision item ("why an evolutionary algorithm rather than
% something simpler?"). Included from main.tex at the end of Section VII via
%   % greedy_baseline.tex
% -----------------------------------------------------------------------------
% Material for the B2 revision item ("why an evolutionary algorithm rather than
% something simpler?"). Included from main.tex at the end of Section VII via
%   \input{greedy_baseline}
%
% Requires (already in main.tex's preamble): algorithm, algpseudocode, booktabs.
% The convergence figure it discusses is main.tex's Fig. 1 (Fig1.pdf), which
% already carries the greedy curve -- this file does not duplicate it.
%
% Every number below is measured under the ADOPTED configuration: Algorithm 1
% with the canonical link-count update, Algorithm 3 pruning-only, no repair on
% the stable grids. Producing code:
%   EA functions/greedy_prune.m      Algorithm 3 (the greedy search itself)
%   EA functions/greedy_codesign.m   the bidirectional variant, kept for reference
%   analysis_perf_bounds.m           Fig. 1 with the greedy curve
%   scratchpad/remeasure_canonical.m the eight-seed comparison and Jaccard data
%   scratchpad/ell_freefall.m        the plateau widths quoted below
%   scratchpad/remeasure_prune2.m    the unstable-plant numbers
%
% NOTE for the camera-ready: Table 2 and Fig. 1 both cover seeds 1, 2 and 10,
% which were chosen because the EA does well on them -- a selection, not a
% random draw. The EA column is a mean over five runs per instance: measured
% run-to-run CV reaches 12.9% on the 7x7 grid, so the earlier single-run version
% of this table was not reproducible (it recorded two "EA wins" that 0 of 5
% repetitions reproduce). Producing script: scratchpad/b2_repeats.m.
% -----------------------------------------------------------------------------

\subsection{A greedy baseline}
\label{sec:greedy}

The natural objection to an evolutionary algorithm is that a simpler local
search might do as well. We test this directly rather than argue it. Algorithm
\ref{alg:greedy} searches the \emph{same} gene $\theta = [\ell, \mathbf{a},
\mathbf{s}]$ as Algorithm \ref{alg:ea}, decodes it through the same map
$K_\mathrm{s}(\theta)$, and scores it with the same cost. Its move set is
restricted to \emph{removals}: nothing once pruned is ever restored, so the
retained support is nested decreasing over the whole run. It starts from the
dense architecture $\theta_0 = [\|K_\mathrm{d}\|_0, \mathbf{1}, \mathbf{1}]$,
which is also the EA's initialization, and cycles three phases, each run to its
own fixed point.

\begin{algorithm}%[t]
\caption{Greedy pruning (monotone backward elimination)}
\label{alg:greedy}
\begin{algorithmic}[1]
\Statex \textbf{Input:}
\Statex \quad Plant matrices $A, B$; dense gain $K_\mathrm{d}$
\Statex \quad Link mutation range $d$ (the same $d$ as Algorithm \ref{alg:ea})
\State $\ell \gets \|K_\mathrm{d}\|_0$, \; $\mathbf{a} \gets \mathbf{1}$, \; $\mathbf{s} \gets \mathbf{1}$
\State $J \gets J_\mathrm{EA}([\ell,\mathbf{a},\mathbf{s}])$
\State \textbf{repeat}
\Statex \quad \textit{Phase 1: link count, decrements only}
\State \quad \textbf{while} some $\delta \in \{-1, -d\}$ improves $J$\textbf{:}
\State \quad \quad $\ell \gets \mathrm{sat}(\ell + \delta^\star)$ for the best such $\delta^\star$
\Statex \quad \textit{Phase 2: actuators, deactivations only}
\State \quad \textbf{while} some $u$ with $a_u = 1$ improves $J$ when set to $0$\textbf{:}
\State \quad \quad $a_{u^\star} \gets 0$ for the best such $u^\star$
\Statex \quad \textit{Phase 3: sensors, deactivations only}
\State \quad \textbf{while} some $j$ with $s_j = 1$ improves $J$ when set to $0$\textbf{:}
\State \quad \quad $s_{j^\star} \gets 0$ for the best such $j^\star$
\State \textbf{until} a full cycle yields no improvement
\State \Return $[\ell, \mathbf{a}, \mathbf{s}]$
\end{algorithmic}
\end{algorithm}

Three choices make the comparison like-for-like rather than favorable to either
side. First, Phase 1 moves $\ell$ only by $-1$ and $-d$: this is within the EA's
own link neighborhood $\delta \sim \mathrm{Unif}\{-d,\dots,d\}$, so the greedy
search is not handed a global scan over $\ell$ that Algorithm \ref{alg:ea}
structurally cannot perform.\footnote{Exhaustively scanning $\ell$ at the greedy
fixed point's masks returns the identical optimum on five of six grid instances;
on the sixth the monotone descent overshoots, stopping at $\ell = 15$ where
$\ell = 17$ was better by $2.2\%$. Forbidding backtracking has a price, and this
is it.} Second, neither search is charged for work it can prove redundant: the
EA reuses its elites' costs instead of recomputing them, and the greedy search
skips a mask bit whose row or column is already empty, since clearing it cannot
change $K$. Both exclusions are exact --- neither alters the sequence of
iterates --- and charging one but not the other would move the shared axis by
about a factor of two. Third, destabilizing candidates are \emph{scored}, at
infinite cost, rather than skipped, so the reported budget is what the search
actually costs.

Both searches are therefore comparable on a single currency: the number of
evaluations of $J_\mathrm{LQR}$, each one discrete Lyapunov solve. Algorithm
\ref{alg:ea} spends $N_p$ of them in the first generation and $N_p - n_e$
thereafter --- $1510$ at $G_{\max} = 150$ --- and Fig.~\ref{fig:perf-conv} plots
the greedy incumbent against that cumulative count.

\begin{table}[t]
\centering
\caption{Final $J_\mathrm{EA}$, and the cost of getting there. The last column
is Algorithm \ref{alg:greedy}'s evaluation count divided by Algorithm
\ref{alg:ea}'s, both after excluding provably redundant work; it crosses $1$
between $N_u = 49$ and $N_u = 225$. Algorithm \ref{alg:greedy} is
deterministic; Algorithm \ref{alg:ea} is not, and is reported as mean $\pm$
standard deviation over five runs where five were run. The $15\times15$ row uses
the length-scaled $p_m = 4/n$ and $G_{\max} = 400$ of Remark \ref{rem:pm}
rather than the $p_m = 0.05$, $G_{\max} = 150$ of the rows above; at those
settings the EA returns $84.7$.}
\label{tab:greedy}
\begin{tabular}{lrrc}
\toprule
Plant, seed & EA & Greedy & Ratio \\
\midrule
$5\!\times\!5$, 1   & $\mathbf{6.68 \pm 0.00}$ & $7.46$ & $0.13\times$ \\
$5\!\times\!5$, 2   & $5.81 \pm 0.12$ & $\mathbf{5.72}$ & $0.27\times$ \\
$5\!\times\!5$, 10  & $6.42 \pm 0.10$ & $\mathbf{6.33}$ & $0.21\times$ \\
$7\!\times\!7$, 1   & $6.35 \pm 0.83$ & $\mathbf{5.15}$ & $0.38\times$ \\
$7\!\times\!7$, 2   & $\mathbf{6.90 \pm 0.21}$ & $7.07$ & $0.45\times$ \\
$7\!\times\!7$, 10  & $8.64 \pm 1.11$ & $\mathbf{7.74}$ & $0.42\times$ \\
IEEE 13-bus, 1      & $\mathbf{2.45}$ & $2.55$ & $0.15\times$ \\
\midrule
$15\!\times\!15$, 1 & $12.73$ & $\mathbf{11.30}$ & $\mathbf{3.25\times}$ \\
\bottomrule
\end{tabular}
\end{table}

\paragraph{What the comparison shows.}
The greedy search reaches a lower cost than the EA on four of the six grid
instances of Fig.~\ref{fig:perf-conv}, using $13$--$45\%$ of the budget; the EA
wins the remaining two and the IEEE 13-bus system. The last row of Table
\ref{tab:greedy} is the exception that the closing paragraph of this section
develops: by $N_u = 225$ the ratio has inverted, and greedy costs $3.25$ times
what the EA does to reach a cost within $13\%$. The EA does reach the greedy architecture on some runs
--- on $5\times5$ seed $2$ its best of five runs equals the greedy cost to four
decimals, with identical actuator and sensor sets --- but not reliably, and
reporting a single run per instance makes that look like a tie when it is the
minimum of a distribution. Two of the six instances have a standard deviation
above $10\%$ of the mean, so single-run comparisons against a deterministic
baseline are not meaningful at this effect size; the repetitions in Table
\ref{tab:greedy} are what the comparison rests on. Neither non-monotonicity nor
the infinite-cost set protects the EA: the greedy search probes only $4$--$51$
infeasible candidates per run and simply never selects one, whereas the EA
probes $146$--$782$ because mutation samples them unconditionally.

\paragraph{The link coordinate needs a gradient.}
The comparison exposed a defect in the encoding of Section \ref{sec:ea}. Under
the decode of Algorithm \ref{alg:ea}, $\Pi_\ell$ selects first and the masks
zero entries afterwards, so $N_c \ll \ell$ in general: at the returned elite,
$\ell$ could be lowered by $85$--$254$ before $J_\mathrm{EA}$ changed in any
digit. Selection therefore gave $\ell$ no gradient at all, and it moved by only
$1$--$10\%$ from its initial value while all sparsity came from the masks. The
remedy costs one line --- after each decode, reset the gene's $\ell$ to the
deepest rank still present in $K_\mathrm{s}(\theta)$. The controller, the cost,
and the image of the decode map are unchanged; only the redundant slack in the
genotype is removed, so the next $-\delta$ mutation deletes links that exist.
Over a wider sweep of eight seeds (sixteen grid instances, both arms driven by
the same random stream) this lowered the final cost on thirteen, tied three,
and worsened none (mean $7.83 \to 6.73$, sign test $p = 2.4\times10^{-4}$), and
brought $\ell$ from a mean of $480$ down to $56$, on fourteen of the sixteen
into the range the greedy search reaches. We report this because it is the
analysis of Section
\ref{sec:convergence}, not the simulations, that makes the defect visible: a
coordinate on which the objective is constant cannot be driven by selection.

\paragraph{Scope of the certified gap.}
The greedy solutions also delimit what Theorem \ref{thm:gap} claims. Its bound
is on $\min_{\mathbf{a}} J_\mathrm{EA}$ at the elite's \emph{current} $\ell$ and
$\mathbf{s}$, and the greedy search varies both. On one run of the $7\times7$
grid the greedy architecture costs less than that bound ($5.145$ against
$5.462$), which is not a contradiction but does quantify the restriction:
relaxing the two
frozen coordinates can occasionally buy more than the entire certified
actuator-only gap. A second caveat is that Algorithm \ref{alg:greedy} never
tests an addition, so its fixed point is not a $1$-flip local optimum in the
sense of Definition \ref{def:loc-opt} and the slack $\Xi$ of \eqref{eq:slack}
need not vanish there; the bound remains valid, but it is not tight at a greedy
output the way it can be at an EA output.

\paragraph{On the unstable plant the repair is what separates them.}
Section \ref{sec:gersh_repair}'s open-loop unstable plant is the one setting in
which the two searches are not interchangeable. Running both with and without
Algorithm \ref{alg:gers-repair}: without it the greedy search reaches $15.42$
against the EA's $18.56 \pm 0.53$; with it the EA reaches $15.66 \pm 0.36$
against the greedy search's $16.74$. The repair improves the EA by $16\%$ and
\emph{degrades} the greedy search by $9\%$. Making infeasible candidates cheap
but finite is useful to a population that can carry a bad commitment and
discard it later, and harmful to a monotone search that cannot: roughly half of
the greedy probes are infeasible there, and a move that looks good only because
the repair rescued it can never be undone. This is the clearest evidence we have
that the repair mechanism and the population are complementary rather than
independent additions.

% The scaling argument that used to sit here now closes Section VII in main.tex,
% immediately after this file is \input.

%
% Requires (already in main.tex's preamble): algorithm, algpseudocode, booktabs.
% The convergence figure it discusses is main.tex's Fig. 1 (Fig1.pdf), which
% already carries the greedy curve -- this file does not duplicate it.
%
% Every number below is measured under the ADOPTED configuration: Algorithm 1
% with the canonical link-count update, Algorithm 3 pruning-only, no repair on
% the stable grids. Producing code:
%   EA functions/greedy_prune.m      Algorithm 3 (the greedy search itself)
%   EA functions/greedy_codesign.m   the bidirectional variant, kept for reference
%   analysis_perf_bounds.m           Fig. 1 with the greedy curve
%   scratchpad/remeasure_canonical.m the eight-seed comparison and Jaccard data
%   scratchpad/ell_freefall.m        the plateau widths quoted below
%   scratchpad/remeasure_prune2.m    the unstable-plant numbers
%
% NOTE for the camera-ready: Table 2 and Fig. 1 both cover seeds 1, 2 and 10,
% which were chosen because the EA does well on them -- a selection, not a
% random draw. The EA column is a mean over five runs per instance: measured
% run-to-run CV reaches 12.9% on the 7x7 grid, so the earlier single-run version
% of this table was not reproducible (it recorded two "EA wins" that 0 of 5
% repetitions reproduce). Producing script: scratchpad/b2_repeats.m.
% -----------------------------------------------------------------------------

\subsection{A greedy baseline}
\label{sec:greedy}

The natural objection to an evolutionary algorithm is that a simpler local
search might do as well. We test this directly rather than argue it. Algorithm
\ref{alg:greedy} searches the \emph{same} gene $\theta = [\ell, \mathbf{a},
\mathbf{s}]$ as Algorithm \ref{alg:ea}, decodes it through the same map
$K_\mathrm{s}(\theta)$, and scores it with the same cost. Its move set is
restricted to \emph{removals}: nothing once pruned is ever restored, so the
retained support is nested decreasing over the whole run. It starts from the
dense architecture $\theta_0 = [\|K_\mathrm{d}\|_0, \mathbf{1}, \mathbf{1}]$,
which is also the EA's initialization, and cycles three phases, each run to its
own fixed point.

\begin{algorithm}%[t]
\caption{Greedy pruning (monotone backward elimination)}
\label{alg:greedy}
\begin{algorithmic}[1]
\Statex \textbf{Input:}
\Statex \quad Plant matrices $A, B$; dense gain $K_\mathrm{d}$
\Statex \quad Link mutation range $d$ (the same $d$ as Algorithm \ref{alg:ea})
\State $\ell \gets \|K_\mathrm{d}\|_0$, \; $\mathbf{a} \gets \mathbf{1}$, \; $\mathbf{s} \gets \mathbf{1}$
\State $J \gets J_\mathrm{EA}([\ell,\mathbf{a},\mathbf{s}])$
\State \textbf{repeat}
\Statex \quad \textit{Phase 1: link count, decrements only}
\State \quad \textbf{while} some $\delta \in \{-1, -d\}$ improves $J$\textbf{:}
\State \quad \quad $\ell \gets \mathrm{sat}(\ell + \delta^\star)$ for the best such $\delta^\star$
\Statex \quad \textit{Phase 2: actuators, deactivations only}
\State \quad \textbf{while} some $u$ with $a_u = 1$ improves $J$ when set to $0$\textbf{:}
\State \quad \quad $a_{u^\star} \gets 0$ for the best such $u^\star$
\Statex \quad \textit{Phase 3: sensors, deactivations only}
\State \quad \textbf{while} some $j$ with $s_j = 1$ improves $J$ when set to $0$\textbf{:}
\State \quad \quad $s_{j^\star} \gets 0$ for the best such $j^\star$
\State \textbf{until} a full cycle yields no improvement
\State \Return $[\ell, \mathbf{a}, \mathbf{s}]$
\end{algorithmic}
\end{algorithm}

Three choices make the comparison like-for-like rather than favorable to either
side. First, Phase 1 moves $\ell$ only by $-1$ and $-d$: this is within the EA's
own link neighborhood $\delta \sim \mathrm{Unif}\{-d,\dots,d\}$, so the greedy
search is not handed a global scan over $\ell$ that Algorithm \ref{alg:ea}
structurally cannot perform.\footnote{Exhaustively scanning $\ell$ at the greedy
fixed point's masks returns the identical optimum on five of six grid instances;
on the sixth the monotone descent overshoots, stopping at $\ell = 15$ where
$\ell = 17$ was better by $2.2\%$. Forbidding backtracking has a price, and this
is it.} Second, neither search is charged for work it can prove redundant: the
EA reuses its elites' costs instead of recomputing them, and the greedy search
skips a mask bit whose row or column is already empty, since clearing it cannot
change $K$. Both exclusions are exact --- neither alters the sequence of
iterates --- and charging one but not the other would move the shared axis by
about a factor of two. Third, destabilizing candidates are \emph{scored}, at
infinite cost, rather than skipped, so the reported budget is what the search
actually costs.

Both searches are therefore comparable on a single currency: the number of
evaluations of $J_\mathrm{LQR}$, each one discrete Lyapunov solve. Algorithm
\ref{alg:ea} spends $N_p$ of them in the first generation and $N_p - n_e$
thereafter --- $1510$ at $G_{\max} = 150$ --- and Fig.~\ref{fig:perf-conv} plots
the greedy incumbent against that cumulative count.

\begin{table}[t]
\centering
\caption{Final $J_\mathrm{EA}$, and the cost of getting there. The last column
is Algorithm \ref{alg:greedy}'s evaluation count divided by Algorithm
\ref{alg:ea}'s, both after excluding provably redundant work; it crosses $1$
between $N_u = 49$ and $N_u = 225$. Algorithm \ref{alg:greedy} is
deterministic; Algorithm \ref{alg:ea} is not, and is reported as mean $\pm$
standard deviation over five runs where five were run. The $15\times15$ row uses
the length-scaled $p_m = 4/n$ and $G_{\max} = 400$ of Remark \ref{rem:pm}
rather than the $p_m = 0.05$, $G_{\max} = 150$ of the rows above; at those
settings the EA returns $84.7$.}
\label{tab:greedy}
\begin{tabular}{lrrc}
\toprule
Plant, seed & EA & Greedy & Ratio \\
\midrule
$5\!\times\!5$, 1   & $\mathbf{6.68 \pm 0.00}$ & $7.46$ & $0.13\times$ \\
$5\!\times\!5$, 2   & $5.81 \pm 0.12$ & $\mathbf{5.72}$ & $0.27\times$ \\
$5\!\times\!5$, 10  & $6.42 \pm 0.10$ & $\mathbf{6.33}$ & $0.21\times$ \\
$7\!\times\!7$, 1   & $6.35 \pm 0.83$ & $\mathbf{5.15}$ & $0.38\times$ \\
$7\!\times\!7$, 2   & $\mathbf{6.90 \pm 0.21}$ & $7.07$ & $0.45\times$ \\
$7\!\times\!7$, 10  & $8.64 \pm 1.11$ & $\mathbf{7.74}$ & $0.42\times$ \\
IEEE 13-bus, 1      & $\mathbf{2.45}$ & $2.55$ & $0.15\times$ \\
\midrule
$15\!\times\!15$, 1 & $12.73$ & $\mathbf{11.30}$ & $\mathbf{3.25\times}$ \\
\bottomrule
\end{tabular}
\end{table}

\paragraph{What the comparison shows.}
The greedy search reaches a lower cost than the EA on four of the six grid
instances of Fig.~\ref{fig:perf-conv}, using $13$--$45\%$ of the budget; the EA
wins the remaining two and the IEEE 13-bus system. The last row of Table
\ref{tab:greedy} is the exception that the closing paragraph of this section
develops: by $N_u = 225$ the ratio has inverted, and greedy costs $3.25$ times
what the EA does to reach a cost within $13\%$. The EA does reach the greedy architecture on some runs
--- on $5\times5$ seed $2$ its best of five runs equals the greedy cost to four
decimals, with identical actuator and sensor sets --- but not reliably, and
reporting a single run per instance makes that look like a tie when it is the
minimum of a distribution. Two of the six instances have a standard deviation
above $10\%$ of the mean, so single-run comparisons against a deterministic
baseline are not meaningful at this effect size; the repetitions in Table
\ref{tab:greedy} are what the comparison rests on. Neither non-monotonicity nor
the infinite-cost set protects the EA: the greedy search probes only $4$--$51$
infeasible candidates per run and simply never selects one, whereas the EA
probes $146$--$782$ because mutation samples them unconditionally.

\paragraph{The link coordinate needs a gradient.}
The comparison exposed a defect in the encoding of Section \ref{sec:ea}. Under
the decode of Algorithm \ref{alg:ea}, $\Pi_\ell$ selects first and the masks
zero entries afterwards, so $N_c \ll \ell$ in general: at the returned elite,
$\ell$ could be lowered by $85$--$254$ before $J_\mathrm{EA}$ changed in any
digit. Selection therefore gave $\ell$ no gradient at all, and it moved by only
$1$--$10\%$ from its initial value while all sparsity came from the masks. The
remedy costs one line --- after each decode, reset the gene's $\ell$ to the
deepest rank still present in $K_\mathrm{s}(\theta)$. The controller, the cost,
and the image of the decode map are unchanged; only the redundant slack in the
genotype is removed, so the next $-\delta$ mutation deletes links that exist.
Over a wider sweep of eight seeds (sixteen grid instances, both arms driven by
the same random stream) this lowered the final cost on thirteen, tied three,
and worsened none (mean $7.83 \to 6.73$, sign test $p = 2.4\times10^{-4}$), and
brought $\ell$ from a mean of $480$ down to $56$, on fourteen of the sixteen
into the range the greedy search reaches. We report this because it is the
analysis of Section
\ref{sec:convergence}, not the simulations, that makes the defect visible: a
coordinate on which the objective is constant cannot be driven by selection.

\paragraph{Scope of the certified gap.}
The greedy solutions also delimit what Theorem \ref{thm:gap} claims. Its bound
is on $\min_{\mathbf{a}} J_\mathrm{EA}$ at the elite's \emph{current} $\ell$ and
$\mathbf{s}$, and the greedy search varies both. On one run of the $7\times7$
grid the greedy architecture costs less than that bound ($5.145$ against
$5.462$), which is not a contradiction but does quantify the restriction:
relaxing the two
frozen coordinates can occasionally buy more than the entire certified
actuator-only gap. A second caveat is that Algorithm \ref{alg:greedy} never
tests an addition, so its fixed point is not a $1$-flip local optimum in the
sense of Definition \ref{def:loc-opt} and the slack $\Xi$ of \eqref{eq:slack}
need not vanish there; the bound remains valid, but it is not tight at a greedy
output the way it can be at an EA output.

\paragraph{On the unstable plant the repair is what separates them.}
Section \ref{sec:gersh_repair}'s open-loop unstable plant is the one setting in
which the two searches are not interchangeable. Running both with and without
Algorithm \ref{alg:gers-repair}: without it the greedy search reaches $15.42$
against the EA's $18.56 \pm 0.53$; with it the EA reaches $15.66 \pm 0.36$
against the greedy search's $16.74$. The repair improves the EA by $16\%$ and
\emph{degrades} the greedy search by $9\%$. Making infeasible candidates cheap
but finite is useful to a population that can carry a bad commitment and
discard it later, and harmful to a monotone search that cannot: roughly half of
the greedy probes are infeasible there, and a move that looks good only because
the repair rescued it can never be undone. This is the clearest evidence we have
that the repair mechanism and the population are complementary rather than
independent additions.

% The scaling argument that used to sit here now closes Section VII in main.tex,
% immediately after this file is \input.

%
% Requires (already in main.tex's preamble): algorithm, algpseudocode, booktabs.
% The convergence figure it discusses is main.tex's Fig. 1 (Fig1.pdf), which
% already carries the greedy curve -- this file does not duplicate it.
%
% Every number below is measured under the ADOPTED configuration: Algorithm 1
% with the canonical link-count update, Algorithm 3 pruning-only, no repair on
% the stable grids. Producing code:
%   EA functions/greedy_prune.m      Algorithm 3 (the greedy search itself)
%   EA functions/greedy_codesign.m   the bidirectional variant, kept for reference
%   analysis_perf_bounds.m           Fig. 1 with the greedy curve
%   scratchpad/remeasure_canonical.m the eight-seed comparison and Jaccard data
%   scratchpad/ell_freefall.m        the plateau widths quoted below
%   scratchpad/remeasure_prune2.m    the unstable-plant numbers
%
% NOTE for the camera-ready: Table 2 and Fig. 1 both cover seeds 1, 2 and 10,
% which were chosen because the EA does well on them -- a selection, not a
% random draw. The EA column is a mean over five runs per instance: measured
% run-to-run CV reaches 12.9% on the 7x7 grid, so the earlier single-run version
% of this table was not reproducible (it recorded two "EA wins" that 0 of 5
% repetitions reproduce). Producing script: scratchpad/b2_repeats.m.
% -----------------------------------------------------------------------------

\subsection{A greedy baseline}
\label{sec:greedy}

The natural objection to an evolutionary algorithm is that a simpler local
search might do as well. We test this directly rather than argue it. Algorithm
\ref{alg:greedy} searches the \emph{same} gene $\theta = [\ell, \mathbf{a},
\mathbf{s}]$ as Algorithm \ref{alg:ea}, decodes it through the same map
$K_\mathrm{s}(\theta)$, and scores it with the same cost. Its move set is
restricted to \emph{removals}: nothing once pruned is ever restored, so the
retained support is nested decreasing over the whole run. It starts from the
dense architecture $\theta_0 = [\|K_\mathrm{d}\|_0, \mathbf{1}, \mathbf{1}]$,
which is also the EA's initialization, and cycles three phases, each run to its
own fixed point.

\begin{algorithm}%[t]
\caption{Greedy pruning (monotone backward elimination)}
\label{alg:greedy}
\begin{algorithmic}[1]
\Statex \textbf{Input:}
\Statex \quad Plant matrices $A, B$; dense gain $K_\mathrm{d}$
\Statex \quad Link mutation range $d$ (the same $d$ as Algorithm \ref{alg:ea})
\State $\ell \gets \|K_\mathrm{d}\|_0$, \; $\mathbf{a} \gets \mathbf{1}$, \; $\mathbf{s} \gets \mathbf{1}$
\State $J \gets J_\mathrm{EA}([\ell,\mathbf{a},\mathbf{s}])$
\State \textbf{repeat}
\Statex \quad \textit{Phase 1: link count, decrements only}
\State \quad \textbf{while} some $\delta \in \{-1, -d\}$ improves $J$\textbf{:}
\State \quad \quad $\ell \gets \mathrm{sat}(\ell + \delta^\star)$ for the best such $\delta^\star$
\Statex \quad \textit{Phase 2: actuators, deactivations only}
\State \quad \textbf{while} some $u$ with $a_u = 1$ improves $J$ when set to $0$\textbf{:}
\State \quad \quad $a_{u^\star} \gets 0$ for the best such $u^\star$
\Statex \quad \textit{Phase 3: sensors, deactivations only}
\State \quad \textbf{while} some $j$ with $s_j = 1$ improves $J$ when set to $0$\textbf{:}
\State \quad \quad $s_{j^\star} \gets 0$ for the best such $j^\star$
\State \textbf{until} a full cycle yields no improvement
\State \Return $[\ell, \mathbf{a}, \mathbf{s}]$
\end{algorithmic}
\end{algorithm}

Three choices make the comparison like-for-like rather than favorable to either
side. First, Phase 1 moves $\ell$ only by $-1$ and $-d$: this is within the EA's
own link neighborhood $\delta \sim \mathrm{Unif}\{-d,\dots,d\}$, so the greedy
search is not handed a global scan over $\ell$ that Algorithm \ref{alg:ea}
structurally cannot perform.\footnote{Exhaustively scanning $\ell$ at the greedy
fixed point's masks returns the identical optimum on five of six grid instances;
on the sixth the monotone descent overshoots, stopping at $\ell = 15$ where
$\ell = 17$ was better by $2.2\%$. Forbidding backtracking has a price, and this
is it.} Second, neither search is charged for work it can prove redundant: the
EA reuses its elites' costs instead of recomputing them, and the greedy search
skips a mask bit whose row or column is already empty, since clearing it cannot
change $K$. Both exclusions are exact --- neither alters the sequence of
iterates --- and charging one but not the other would move the shared axis by
about a factor of two. Third, destabilizing candidates are \emph{scored}, at
infinite cost, rather than skipped, so the reported budget is what the search
actually costs.

Both searches are therefore comparable on a single currency: the number of
evaluations of $J_\mathrm{LQR}$, each one discrete Lyapunov solve. Algorithm
\ref{alg:ea} spends $N_p$ of them in the first generation and $N_p - n_e$
thereafter --- $1510$ at $G_{\max} = 150$ --- and Fig.~\ref{fig:perf-conv} plots
the greedy incumbent against that cumulative count.

\begin{table}[t]
\centering
\caption{Final $J_\mathrm{EA}$, and the cost of getting there. The last column
is Algorithm \ref{alg:greedy}'s evaluation count divided by Algorithm
\ref{alg:ea}'s, both after excluding provably redundant work; it crosses $1$
between $N_u = 49$ and $N_u = 225$. Algorithm \ref{alg:greedy} is
deterministic; Algorithm \ref{alg:ea} is not, and is reported as mean $\pm$
standard deviation over five runs where five were run. The $15\times15$ row uses
the length-scaled $p_m = 4/n$ and $G_{\max} = 400$ of Remark \ref{rem:pm}
rather than the $p_m = 0.05$, $G_{\max} = 150$ of the rows above; at those
settings the EA returns $84.7$.}
\label{tab:greedy}
\begin{tabular}{lrrc}
\toprule
Plant, seed & EA & Greedy & Ratio \\
\midrule
$5\!\times\!5$, 1   & $\mathbf{6.68 \pm 0.00}$ & $7.46$ & $0.13\times$ \\
$5\!\times\!5$, 2   & $5.81 \pm 0.12$ & $\mathbf{5.72}$ & $0.27\times$ \\
$5\!\times\!5$, 10  & $6.42 \pm 0.10$ & $\mathbf{6.33}$ & $0.21\times$ \\
$7\!\times\!7$, 1   & $6.35 \pm 0.83$ & $\mathbf{5.15}$ & $0.38\times$ \\
$7\!\times\!7$, 2   & $\mathbf{6.90 \pm 0.21}$ & $7.07$ & $0.45\times$ \\
$7\!\times\!7$, 10  & $8.64 \pm 1.11$ & $\mathbf{7.74}$ & $0.42\times$ \\
IEEE 13-bus, 1      & $\mathbf{2.45}$ & $2.55$ & $0.15\times$ \\
\midrule
$15\!\times\!15$, 1 & $12.73$ & $\mathbf{11.30}$ & $\mathbf{3.25\times}$ \\
\bottomrule
\end{tabular}
\end{table}

\paragraph{What the comparison shows.}
The greedy search reaches a lower cost than the EA on four of the six grid
instances of Fig.~\ref{fig:perf-conv}, using $13$--$45\%$ of the budget; the EA
wins the remaining two and the IEEE 13-bus system. The last row of Table
\ref{tab:greedy} is the exception that the closing paragraph of this section
develops: by $N_u = 225$ the ratio has inverted, and greedy costs $3.25$ times
what the EA does to reach a cost within $13\%$. The EA does reach the greedy architecture on some runs
--- on $5\times5$ seed $2$ its best of five runs equals the greedy cost to four
decimals, with identical actuator and sensor sets --- but not reliably, and
reporting a single run per instance makes that look like a tie when it is the
minimum of a distribution. Two of the six instances have a standard deviation
above $10\%$ of the mean, so single-run comparisons against a deterministic
baseline are not meaningful at this effect size; the repetitions in Table
\ref{tab:greedy} are what the comparison rests on. Neither non-monotonicity nor
the infinite-cost set protects the EA: the greedy search probes only $4$--$51$
infeasible candidates per run and simply never selects one, whereas the EA
probes $146$--$782$ because mutation samples them unconditionally.

\paragraph{The link coordinate needs a gradient.}
The comparison exposed a defect in the encoding of Section \ref{sec:ea}. Under
the decode of Algorithm \ref{alg:ea}, $\Pi_\ell$ selects first and the masks
zero entries afterwards, so $N_c \ll \ell$ in general: at the returned elite,
$\ell$ could be lowered by $85$--$254$ before $J_\mathrm{EA}$ changed in any
digit. Selection therefore gave $\ell$ no gradient at all, and it moved by only
$1$--$10\%$ from its initial value while all sparsity came from the masks. The
remedy costs one line --- after each decode, reset the gene's $\ell$ to the
deepest rank still present in $K_\mathrm{s}(\theta)$. The controller, the cost,
and the image of the decode map are unchanged; only the redundant slack in the
genotype is removed, so the next $-\delta$ mutation deletes links that exist.
Over a wider sweep of eight seeds (sixteen grid instances, both arms driven by
the same random stream) this lowered the final cost on thirteen, tied three,
and worsened none (mean $7.83 \to 6.73$, sign test $p = 2.4\times10^{-4}$), and
brought $\ell$ from a mean of $480$ down to $56$, on fourteen of the sixteen
into the range the greedy search reaches. We report this because it is the
analysis of Section
\ref{sec:convergence}, not the simulations, that makes the defect visible: a
coordinate on which the objective is constant cannot be driven by selection.

\paragraph{Scope of the certified gap.}
The greedy solutions also delimit what Theorem \ref{thm:gap} claims. Its bound
is on $\min_{\mathbf{a}} J_\mathrm{EA}$ at the elite's \emph{current} $\ell$ and
$\mathbf{s}$, and the greedy search varies both. On one run of the $7\times7$
grid the greedy architecture costs less than that bound ($5.145$ against
$5.462$), which is not a contradiction but does quantify the restriction:
relaxing the two
frozen coordinates can occasionally buy more than the entire certified
actuator-only gap. A second caveat is that Algorithm \ref{alg:greedy} never
tests an addition, so its fixed point is not a $1$-flip local optimum in the
sense of Definition \ref{def:loc-opt} and the slack $\Xi$ of \eqref{eq:slack}
need not vanish there; the bound remains valid, but it is not tight at a greedy
output the way it can be at an EA output.

\paragraph{On the unstable plant the repair is what separates them.}
Section \ref{sec:gersh_repair}'s open-loop unstable plant is the one setting in
which the two searches are not interchangeable. Running both with and without
Algorithm \ref{alg:gers-repair}: without it the greedy search reaches $15.42$
against the EA's $18.56 \pm 0.53$; with it the EA reaches $15.66 \pm 0.36$
against the greedy search's $16.74$. The repair improves the EA by $16\%$ and
\emph{degrades} the greedy search by $9\%$. Making infeasible candidates cheap
but finite is useful to a population that can carry a bad commitment and
discard it later, and harmful to a monotone search that cannot: roughly half of
the greedy probes are infeasible there, and a move that looks good only because
the repair rescued it can never be undone. This is the clearest evidence we have
that the repair mechanism and the population are complementary rather than
independent additions.

% The scaling argument that used to sit here now closes Section VII in main.tex,
% immediately after this file is \input.

%
% Requires (already in main.tex's preamble): algorithm, algpseudocode, booktabs.
% The convergence figure it discusses is main.tex's Fig. 1 (Fig1.pdf), which
% already carries the greedy curve -- this file does not duplicate it.
%
% Every number below is measured under the ADOPTED configuration: Algorithm 1
% with the canonical link-count update, Algorithm 3 pruning-only, no repair on
% the stable grids. Producing code:
%   EA functions/greedy_prune.m      Algorithm 3 (the greedy search itself)
%   EA functions/greedy_codesign.m   the bidirectional variant, kept for reference
%   analysis_perf_bounds.m           Fig. 1 with the greedy curve
%   scratchpad/remeasure_canonical.m the eight-seed comparison and Jaccard data
%   scratchpad/ell_freefall.m        the plateau widths quoted below
%   scratchpad/remeasure_prune2.m    the unstable-plant numbers
%
% NOTE for the camera-ready: Table 2 and Fig. 1 both cover seeds 1, 2 and 10,
% which were chosen because the EA does well on them -- a selection, not a
% random draw. The EA column is a mean over five runs per instance: measured
% run-to-run CV reaches 12.9% on the 7x7 grid, so the earlier single-run version
% of this table was not reproducible (it recorded two "EA wins" that 0 of 5
% repetitions reproduce). Producing script: scratchpad/b2_repeats.m.
% -----------------------------------------------------------------------------

\subsection{A greedy baseline}
\label{sec:greedy}

The natural objection to an evolutionary algorithm is that a simpler local
search might do as well. We test this directly rather than argue it. Algorithm
\ref{alg:greedy} searches the \emph{same} gene $\theta = [\ell, \mathbf{a},
\mathbf{s}]$ as Algorithm \ref{alg:ea}, decodes it through the same map
$K_\mathrm{s}(\theta)$, and scores it with the same cost. Its move set is
restricted to \emph{removals}: nothing once pruned is ever restored, so the
retained support is nested decreasing over the whole run. It starts from the
dense architecture $\theta_0 = [\|K_\mathrm{d}\|_0, \mathbf{1}, \mathbf{1}]$,
which is also the EA's initialization, and cycles three phases, each run to its
own fixed point.

\begin{algorithm}%[t]
\caption{Greedy pruning (monotone backward elimination)}
\label{alg:greedy}
\begin{algorithmic}[1]
\Statex \textbf{Input:}
\Statex \quad Plant matrices $A, B$; dense gain $K_\mathrm{d}$
\Statex \quad Link mutation range $d$ (the same $d$ as Algorithm \ref{alg:ea})
\State $\ell \gets \|K_\mathrm{d}\|_0$, \; $\mathbf{a} \gets \mathbf{1}$, \; $\mathbf{s} \gets \mathbf{1}$
\State $J \gets J_\mathrm{EA}([\ell,\mathbf{a},\mathbf{s}])$
\State \textbf{repeat}
\Statex \quad \textit{Phase 1: link count, decrements only}
\State \quad \textbf{while} some $\delta \in \{-1, -d\}$ improves $J$\textbf{:}
\State \quad \quad $\ell \gets \mathrm{sat}(\ell + \delta^\star)$ for the best such $\delta^\star$
\Statex \quad \textit{Phase 2: actuators, deactivations only}
\State \quad \textbf{while} some $u$ with $a_u = 1$ improves $J$ when set to $0$\textbf{:}
\State \quad \quad $a_{u^\star} \gets 0$ for the best such $u^\star$
\Statex \quad \textit{Phase 3: sensors, deactivations only}
\State \quad \textbf{while} some $j$ with $s_j = 1$ improves $J$ when set to $0$\textbf{:}
\State \quad \quad $s_{j^\star} \gets 0$ for the best such $j^\star$
\State \textbf{until} a full cycle yields no improvement
\State \Return $[\ell, \mathbf{a}, \mathbf{s}]$
\end{algorithmic}
\end{algorithm}

Three choices make the comparison like-for-like rather than favorable to either
side. First, Phase 1 moves $\ell$ only by $-1$ and $-d$: this is within the EA's
own link neighborhood $\delta \sim \mathrm{Unif}\{-d,\dots,d\}$, so the greedy
search is not handed a global scan over $\ell$ that Algorithm \ref{alg:ea}
structurally cannot perform.\footnote{Exhaustively scanning $\ell$ at the greedy
fixed point's masks returns the identical optimum on five of six grid instances;
on the sixth the monotone descent overshoots, stopping at $\ell = 15$ where
$\ell = 17$ was better by $2.2\%$. Forbidding backtracking has a price, and this
is it.} Second, neither search is charged for work it can prove redundant: the
EA reuses its elites' costs instead of recomputing them, and the greedy search
skips a mask bit whose row or column is already empty, since clearing it cannot
change $K$. Both exclusions are exact --- neither alters the sequence of
iterates --- and charging one but not the other would move the shared axis by
about a factor of two. Third, destabilizing candidates are \emph{scored}, at
infinite cost, rather than skipped, so the reported budget is what the search
actually costs.

Both searches are therefore comparable on a single currency: the number of
evaluations of $J_\mathrm{LQR}$, each one discrete Lyapunov solve. Algorithm
\ref{alg:ea} spends $N_p$ of them in the first generation and $N_p - n_e$
thereafter --- $1510$ at $G_{\max} = 150$ --- and Fig.~\ref{fig:perf-conv} plots
the greedy incumbent against that cumulative count.

\begin{table}[t]
\centering
\caption{Final $J_\mathrm{EA}$, and the cost of getting there. The last column
is Algorithm \ref{alg:greedy}'s evaluation count divided by Algorithm
\ref{alg:ea}'s, both after excluding provably redundant work; it crosses $1$
between $N_u = 49$ and $N_u = 225$. Algorithm \ref{alg:greedy} is
deterministic; Algorithm \ref{alg:ea} is not, and is reported as mean $\pm$
standard deviation over five runs where five were run. The $15\times15$ row uses
the length-scaled $p_m = 4/n$ and $G_{\max} = 400$ of Remark \ref{rem:pm}
rather than the $p_m = 0.05$, $G_{\max} = 150$ of the rows above; at those
settings the EA returns $84.7$.}
\label{tab:greedy}
\begin{tabular}{lrrc}
\toprule
Plant, seed & EA & Greedy & Ratio \\
\midrule
$5\!\times\!5$, 1   & $\mathbf{6.68 \pm 0.00}$ & $7.46$ & $0.13\times$ \\
$5\!\times\!5$, 2   & $5.81 \pm 0.12$ & $\mathbf{5.72}$ & $0.27\times$ \\
$5\!\times\!5$, 10  & $6.42 \pm 0.10$ & $\mathbf{6.33}$ & $0.21\times$ \\
$7\!\times\!7$, 1   & $6.35 \pm 0.83$ & $\mathbf{5.15}$ & $0.38\times$ \\
$7\!\times\!7$, 2   & $\mathbf{6.90 \pm 0.21}$ & $7.07$ & $0.45\times$ \\
$7\!\times\!7$, 10  & $8.64 \pm 1.11$ & $\mathbf{7.74}$ & $0.42\times$ \\
IEEE 13-bus, 1      & $\mathbf{2.45}$ & $2.55$ & $0.15\times$ \\
\midrule
$15\!\times\!15$, 1 & $12.73$ & $\mathbf{11.30}$ & $\mathbf{3.25\times}$ \\
\bottomrule
\end{tabular}
\end{table}

\paragraph{What the comparison shows.}
The greedy search reaches a lower cost than the EA on four of the six grid
instances of Fig.~\ref{fig:perf-conv}, using $13$--$45\%$ of the budget; the EA
wins the remaining two and the IEEE 13-bus system. The last row of Table
\ref{tab:greedy} is the exception that the closing paragraph of this section
develops: by $N_u = 225$ the ratio has inverted, and greedy costs $3.25$ times
what the EA does to reach a cost within $13\%$. The EA does reach the greedy architecture on some runs
--- on $5\times5$ seed $2$ its best of five runs equals the greedy cost to four
decimals, with identical actuator and sensor sets --- but not reliably, and
reporting a single run per instance makes that look like a tie when it is the
minimum of a distribution. Two of the six instances have a standard deviation
above $10\%$ of the mean, so single-run comparisons against a deterministic
baseline are not meaningful at this effect size; the repetitions in Table
\ref{tab:greedy} are what the comparison rests on. Neither non-monotonicity nor
the infinite-cost set protects the EA: the greedy search probes only $4$--$51$
infeasible candidates per run and simply never selects one, whereas the EA
probes $146$--$782$ because mutation samples them unconditionally.

\paragraph{The link coordinate needs a gradient.}
The comparison exposed a defect in the encoding of Section \ref{sec:ea}. Under
the decode of Algorithm \ref{alg:ea}, $\Pi_\ell$ selects first and the masks
zero entries afterwards, so $N_c \ll \ell$ in general: at the returned elite,
$\ell$ could be lowered by $85$--$254$ before $J_\mathrm{EA}$ changed in any
digit. Selection therefore gave $\ell$ no gradient at all, and it moved by only
$1$--$10\%$ from its initial value while all sparsity came from the masks. The
remedy costs one line --- after each decode, reset the gene's $\ell$ to the
deepest rank still present in $K_\mathrm{s}(\theta)$. The controller, the cost,
and the image of the decode map are unchanged; only the redundant slack in the
genotype is removed, so the next $-\delta$ mutation deletes links that exist.
Over a wider sweep of eight seeds (sixteen grid instances, both arms driven by
the same random stream) this lowered the final cost on thirteen, tied three,
and worsened none (mean $7.83 \to 6.73$, sign test $p = 2.4\times10^{-4}$), and
brought $\ell$ from a mean of $480$ down to $56$, on fourteen of the sixteen
into the range the greedy search reaches. We report this because it is the
analysis of Section
\ref{sec:convergence}, not the simulations, that makes the defect visible: a
coordinate on which the objective is constant cannot be driven by selection.

\paragraph{Scope of the certified gap.}
The greedy solutions also delimit what Theorem \ref{thm:gap} claims. Its bound
is on $\min_{\mathbf{a}} J_\mathrm{EA}$ at the elite's \emph{current} $\ell$ and
$\mathbf{s}$, and the greedy search varies both. On one run of the $7\times7$
grid the greedy architecture costs less than that bound ($5.145$ against
$5.462$), which is not a contradiction but does quantify the restriction:
relaxing the two
frozen coordinates can occasionally buy more than the entire certified
actuator-only gap. A second caveat is that Algorithm \ref{alg:greedy} never
tests an addition, so its fixed point is not a $1$-flip local optimum in the
sense of Definition \ref{def:loc-opt} and the slack $\Xi$ of \eqref{eq:slack}
need not vanish there; the bound remains valid, but it is not tight at a greedy
output the way it can be at an EA output.

\paragraph{On the unstable plant the repair is what separates them.}
Section \ref{sec:gersh_repair}'s open-loop unstable plant is the one setting in
which the two searches are not interchangeable. Running both with and without
Algorithm \ref{alg:gers-repair}: without it the greedy search reaches $15.42$
against the EA's $18.56 \pm 0.53$; with it the EA reaches $15.66 \pm 0.36$
against the greedy search's $16.74$. The repair improves the EA by $16\%$ and
\emph{degrades} the greedy search by $9\%$. Making infeasible candidates cheap
but finite is useful to a population that can carry a bad commitment and
discard it later, and harmful to a monotone search that cannot: roughly half of
the greedy probes are infeasible there, and a move that looks good only because
the repair rescued it can never be undone. This is the clearest evidence we have
that the repair mechanism and the population are complementary rather than
independent additions.

% The scaling argument that used to sit here now closes Section VII in main.tex,
% immediately after this file is \input.
